\begin{table}[htbp]
  \centering
  \caption{Testing whether the firm-size wage premium changed over time}
  \label{tab:firm-size-premium-trend-test}
  \small
  \begin{tabular}{lcccc}
    \toprule
    Firm size & Annual trend & S.E. & $p$-value & Implied 2008--2025 change \\
    \midrule
    2-3 & 0.0024 & (0.0029) & 0.409 & 4.3\% \\
    4-5 & 0.0005 & (0.0039) & 0.909 & 0.8\% \\
    6-10 & 0.0014 & (0.0036) & 0.693 & 2.5\% \\
    11-19 & 0.0022 & (0.0043) & 0.610 & 3.9\% \\
    20-30 & 0.0023 & (0.0042) & 0.595 & 3.9\% \\
    31-50 & 0.0029 & (0.0046) & 0.533 & 5.1\% \\
    51-100 & 0.0037 & (0.0042) & 0.393 & 6.5\% \\
    101+ & 0.0054 & (0.0047) & 0.268 & 9.7\% \\
    \bottomrule
  \end{tabular}
  \vspace{0.3em}
  \begin{minipage}{0.95\textwidth}
  \footnotesize
  Notes: The annual trend is the coefficient on the interaction between firm-size category and a linear year trend, with solo workers as the omitted category. The specification controls for gender, age, age squared, education, formality, and sector-year fixed effects, and uses GEIH expansion weights. Standard errors are clustered by sector. The $p$-value corresponds to the two-sided test that the trend differs from zero. The final column reports $100\times[\exp(17\hat{\delta})-1]$, the implied change in the firm-size wage ratio between 2008 and 2025.
  \end{minipage}
\end{table}
